\section{Introduction}
% While Large Language Models have become common goods at this point \cite{}, Video Language or Video Generative Models have been falling behind in terms of improvement \cite{}. One of their main limitations is that they can generate infeasible trajectories. In other words, videos that do not follow physical laws. This might be undesirable behavior for generating content, but it is completely catastrophic if one wants to train robotics on this generated data. Autonomous robots need to be trained; training in the real world is either expensive or infeasible (cannot let an untrained autonomous car drive into people to learn not to do that!). Therefore, good feasible simulated data is needed to make faster progress in this domain. However, the difficulty lies in that if the data is only slightly off, the robot will learn actions that it cannot execute in the real world. Grounding VGM in true world dynamics is needed.

% TODO: \hl{connect intro to World-Models, DINO and the GVP-WM} 


% Therefore, it is needed to create methods that are grounded

% plan: finish the intro within 0.5 page. 
% - Passage 1: bg
% - Passage 2: our method
% - Passage 3: our contribution

% before compressing: 
% Large video generation models have demonstrated strong zero-shot capabilities in synthesizing realistic and temporally coherent videos, while also generating expert-like video plans conditioned on textual instructions. However, these generated videos often lack physical consistency with the real world, especially in out-of-distribution (OOD) environments, where they may exhibit physical or temporal artifacts such as object teleportation and motion blur.

% To improve the executability of video plans, early works typically treated generated videos as visual subgoals for Model Predictive Control (MPC) or hierarchical reinforcement learning. However, these methods implicitly assume that video plans are always physically feasible, while overlooking their potential physical inconsistencies. More recent approaches introduce goal-conditioned exploration, allowing policies to deviate from erroneous video guidance to mitigate planning errors, but they still rely on additional environment interaction and policy training. To avoid such additional training costs, the original GVP-WM [yixin here: reference to replace the subject] framework proposes a purely test-time planning approach that formulates video grounding as a latent trajectory optimization problem within a pretrained world model. Specifically, it employs an action-conditioned world model as an internal physics simulator and jointly optimizes latent states and actions during inference, thereby projecting physically inconsistent video plans onto dynamically feasible trajectories without requiring additional policy training. However, GVP-WM relies on deterministic hard collocation constraints based on the Augmented Lagrangian Method (ALM), leading to slow test-time optimization. Moreover, its deterministic world model trained with mean squared error (MSE) is brittle under low-quality OOD video guidance and prone to error accumulation in long-horizon planning, ultimately resulting in planning failure.

% In this work, we reproduce GVP-WM [yixin here: reference to replace the subject] from scratch and build the entire system pipeline without access to the original full implementation. During reproduction, we also communicated with the original authors regarding several implementation details. Building upon this reproduction, we further reformulate the original trajectory optimization framework based on deterministic hard constraints into a soft optimization paradigm driven by probabilistic priors. Specifically, we introduce a Score-Based Feasibility Prior by training a conditional denoising network via score matching on expert trajectories to learn the conditional distribution of physically feasible state transitions. We then incorporate the learned prior into the optimization objective as a soft constraint, enabling the model to maintain optimization efficiency while exhibiting stronger robustness under low-quality OOD video guidance. Furthermore, to mitigate error accumulation in long-horizon planning, we introduce a Diagnostic Residual Model that dynamically detects whether latent trajectories gradually drift away from the feasible physical manifold by quantifying the discrepancy between optimized trajectories and predictions from the deterministic world model.

Large video generation models have demonstrated strong zero-shot capabilities in synthesizing realistic and temporally coherent videos while generating expert-like video plans conditioned on textual instructions. However, these generated videos often lack physical consistency with the real world, especially in out-of-distribution (OOD) environments, where they may exhibit physical or temporal artifacts such as object teleportation and motion blur.

To improve the executability of video plans, early works typically treated generated videos as visual subgoals for Model Predictive Control (MPC) or hierarchical reinforcement learning. However, these methods implicitly assume that the generated video plans are always physically feasible, overlooking the physical inconsistencies already present in the videos themselves. More recent approaches alleviate planning errors through goal-conditioned exploration, but still rely on additional environment interaction and policy training. GVP-WM instead proposes a purely test-time planning framework that projects physically inconsistent video plans onto dynamically feasible trajectories through latent trajectory optimization in a pretrained world model. However, its deterministic hard constraints based on the Augmented Lagrangian Method (ALM) lead to slow test-time optimization, while its MSE-trained deterministic world model remains brittle under low-quality OOD video guidance and prone to long-horizon error accumulation.

In this work, we reproduce GVP-WM from scratch and build the entire system pipeline without access to the original full implementation. During reproduction, we also communicated with the original authors regarding several implementation details. Building upon this reproduction, we further reformulate the original deterministic hard-constrained optimization into a soft optimization paradigm based on probabilistic priors. Specifically, we introduce a Score-Based Feasibility Prior as a learned soft constraint to improve robustness under low-quality OOD video guidance. Furthermore, we design a Diagnostic Residual Model to dynamically detect latent trajectory drift from the feasible physical manifold during long-horizon planning.


Our main contributions are as follows:
\begin{itemize}
    \item one
    \item two
\end{itemize}

